% called by main.tex
%
\chapter{Introducción}
\label{ch::introduccion}

Este Trabajo Fin de Máster estudia un problema de predicción financiera de muy corto
plazo en un entorno poco habitual: los mercados de predicción sobre el precio de Bitcoin
de la plataforma Polymarket. El objetivo no es un sistema de trading en producción, sino
una \emph{metodología} rigurosa para responder, con honestidad estadística, a una
pregunta concreta: ¿contiene el libro de órdenes de estos mercados señales que permitan
anticipar el movimiento del precio en los próximos segundos y que, además, sobrevivan a
los costes y a la latencia reales de operar?

Como el lector puede no estar familiarizado con el mundo de las criptomonedas ni con la
microestructura de mercado, este capítulo introduce todos los conceptos desde cero antes
de plantear el problema y los objetivos.

\section{Introducción y definición del problema}

Esta sección define el problema desde cero. Antes de plantear la pregunta de
investigación se introducen los conceptos imprescindibles —mercado de predicción,
contrato binario, libro de órdenes y señales de microestructura—, de modo que el documento
resulte autocontenido también para un lector ajeno al mundo de las criptomonedas. Con ese
vocabulario fijado, se enuncia el problema concreto que se aborda.

\subsection{Qué es un mercado de predicción}

Un \emph{mercado de predicción} es un mercado en el que se compran y venden contratos
cuyo valor depende del resultado de un acontecimiento futuro. El caso más sencillo es el
\emph{contrato binario}: un contrato que pagará una unidad monetaria si el acontecimiento
ocurre y cero si no ocurre. Por construcción, su precio se mueve entre 0 y 1, y puede
leerse directamente como la \emph{probabilidad implícita} que el mercado asigna al
suceso. Por ejemplo, un contrato ``el precio de Bitcoin estará por encima de un umbral a
una hora determinada'' que cotiza a 0,60 indica que el mercado estima una probabilidad
cercana al 60\,\% de que eso ocurra. Cuando el acontecimiento se resuelve, el contrato
vale 1 o 0 y las posiciones se liquidan. La literatura clásica ha mostrado que estos
mercados agregan información dispersa y producen previsiones razonablemente
calibradas~\cite{wolfers2004prediction}.

\subsection{Qué es Polymarket}

Polymarket es una plataforma de mercados de predicción descentralizada. A diferencia de
una casa de apuestas, donde el operador fija las cuotas, en Polymarket son los propios
participantes quienes ponen precios mediante un \emph{libro de órdenes} central
(\emph{Central Limit Order Book}, CLOB): el emparejamiento
de compras y ventas ocurre fuera de la cadena de bloques (de forma rápida) y la
liquidación final del dinero se registra en cadena (sobre una \emph{blockchain})
\cite{polymarket_trading}. En los mercados sobre el precio de Bitcoin, cada contrato
binario toma valores entre 0 y 1 y se interpreta como una
probabilidad~\cite{polymarket_prices}. Esto convierte cada mercado en un pequeño
\emph{mercado electrónico} (\emph{exchange}) con su propia dinámica de oferta y demanda.

Que la liquidación viva en una cadena de bloques tiene una consecuencia directa para un
trabajo de datos: nada se borra. Cada operación liquidada queda registrada de forma
permanente y pública; si una posición se deshace, no se elimina el apunte original, sino
que se añade otro en sentido contrario. Esa propiedad es la que hace posible la auditoría
del registro público de operaciones del Anexo~D —reproducir las cifras del trabajo desde
un dato que nadie puede reescribir— y la que garantiza que el flujo de operaciones
capturado no depende de la buena voluntad del operador de la plataforma. El libro de
órdenes en sí, en cambio, vive fuera de cadena por velocidad, y por eso su captura
continua (Capítulo~\ref{ch::metodos}) es aportación de este trabajo: es la parte del
sistema que no deja huella permanente.

\subsection{El libro de órdenes y el vocabulario de la microestructura}

El \emph{libro de órdenes} (en inglés \emph{limit order book}, LOB) es la lista, ordenada
por precio, de todas las órdenes de compra y de venta que están a la espera de ejecutarse.
Sus elementos esenciales son:

\begin{itemize}
  \item \textbf{Bid}: el precio más alto al que alguien está dispuesto a comprar.
  \item \textbf{Ask}: el precio más bajo al que alguien está dispuesto a vender. Siempre
  es mayor o igual que el \emph{bid}.
  \item \textbf{Spread}: la diferencia entre el \emph{ask} y el \emph{bid}. Es el primer
  coste de operar: quien compra al \emph{ask} y vende al \emph{bid} pierde el \emph{spread}.
  \item \textbf{Profundidad}: el volumen disponible en cada nivel de precio. A más
  profundidad, más se puede operar sin mover el precio.
  \item \textbf{Precio medio} (\emph{mid}): el punto medio entre \emph{bid} y \emph{ask};
  una referencia ``limpia'' del precio sin el ruido del \emph{spread}.
  \item \textbf{Tick}: la variación mínima de precio que admite el mercado. Es la unidad
  natural en la que mediremos todos los resultados de este trabajo.
  \item \textbf{Taker}: quien cruza el \emph{spread} con una orden inmediata contra el
  libro. Obtiene ejecución al instante y, en los mercados de criptomonedas de 2026, paga
  la comisión.
  \item \textbf{Maker} (creador de mercado): quien publica órdenes límite y espera a que
  otro las cruce. Provee liquidez, cobra el \emph{spread} y una parte de las comisiones
  de los \emph{takers} (\emph{rebate}), y a cambio asume el riesgo de quedarse con
  inventario cuando el mercado se mueve en su contra. La distinción \emph{maker}/\emph{taker}
  vertebra la segunda mitad de este trabajo.
\end{itemize}

La Tabla~\ref{tab:lob} muestra, a modo de ejemplo ilustrativo, el aspecto de un libro de
órdenes con tres niveles a cada lado.

\begin{table}[h]
\centering
\caption{Ejemplo ilustrativo de un libro de órdenes (valores inventados con fines
explicativos, no datos reales). El mejor \emph{bid} es 0,52 y el mejor \emph{ask}, 0,54; por
tanto el \emph{spread} es de 0,02 y el precio medio, 0,53. Los niveles más alejados del
centro tienen, típicamente, más volumen acumulado.}
\label{tab:lob}
\begin{tabular}{rr|rr}
\toprule
\multicolumn{2}{c|}{\textbf{Compra (\emph{bids})}} & \multicolumn{2}{c}{\textbf{Venta (\emph{asks})}} \\
Precio & Tamaño & Precio & Tamaño \\
\midrule
0,52 & 120 & 0,54 & \phantom{0}90 \\
0,51 & 200 & 0,55 & 150 \\
0,50 & 340 & 0,56 & 220 \\
\bottomrule
\end{tabular}
\end{table}

La \emph{microestructura de mercado} es la rama que estudia cómo se forman los precios a
partir de estos mecanismos —órdenes, \emph{spread}, profundidad y flujo de
operaciones— a escalas de tiempo muy finas. En este contexto, se entiende por
\emph{señales de microestructura} las variables derivadas del propio libro de órdenes (por
ejemplo, el desbalance entre compra y venta, la presión del microprice o los cambios
recientes del \emph{spread}) que podrían contener información anticipatoria sobre el
movimiento inmediato del precio. La Figura~\ref{fig:mercado} resume el entorno: el precio del
contrato se forma en su propio libro de órdenes, pero está acoplado a referencias externas
del mercado de Bitcoin (mercado al contado, futuros perpetuos y un oráculo de precios), ya
que, en última instancia, el contrato se refiere al precio de Bitcoin.

\begin{figure}[h]
\centering
\begin{tikzpicture}[
  node distance=8mm,
  box/.style={draw, rounded corners, align=center, minimum height=9mm, inner sep=4pt, font=\small},
  ext/.style={box, fill=blue!6},
  core/.style={box, fill=orange!12, thick},
  tgt/.style={box, fill=green!10},
  >=latex
]
\node[ext] (spot) {Mercado al contado\\(spot BTC)};
\node[ext, below=of spot] (perp) {Futuros perpetuos\\(perp BTC)};
\node[ext, below=of perp] (oracle) {Oráculo de precios};
\node[core, right=22mm of perp] (lob) {Libro de órdenes\\del contrato\\(\emph{bid}, \emph{ask}, \emph{spread})};
\node[tgt, right=22mm of lob] (mid) {Precio del contrato\\$m_t \in [0,1]$};
\node[box, below=14mm of mid, fill=gray!8] (pred) {Pregunta: ¿hacia dónde\\se mueve $m_t$ en los\\próximos segundos?};
\draw[->] (spot) -- (lob);
\draw[->] (perp) -- (lob);
\draw[->] (oracle) -- (lob);
\draw[->] (lob) -- (mid);
\draw[->] (mid) -- (pred);
\end{tikzpicture}
\caption{El entorno del problema. El precio del contrato binario se forma en su libro de
órdenes y está acoplado a referencias externas del mercado de Bitcoin. El objetivo es
anticipar su movimiento de muy corto plazo.}
\label{fig:mercado}
\end{figure}

\subsection{Definición del problema}

El precio del contrato no es el precio de Bitcoin: es una función no lineal y dependiente
del tiempo de la probabilidad de que el suceso se resuelva afirmativamente. Sin embargo,
está acoplado a las referencias externas observables que se acaban de citar. Un modelo
útil debe, por tanto, aprender simultáneamente dos elementos: la dinámica interna del libro de
órdenes y la semántica del contrato respecto a Bitcoin.

El problema que se aborda es de microestructura y de muy corto plazo: dado el estado del
mercado en un instante $t$ (libro de órdenes y referencias externas), predecir el
movimiento del precio medio del contrato en una ventana de pocas decenas de segundos. A
diferencia de la mayor parte de la literatura sobre mercados de predicción, que estudia
cómo se calibra el precio frente a la resolución final del evento, aquí el interés está en
los movimientos locales del precio, a escala de segundos.

Conviene fijar desde el principio un matiz que vertebra todo el trabajo: una buena
capacidad de \emph{predecir la dirección} no implica una estrategia rentable. Esta
distinción está bien documentada en la literatura de microestructura y de aprendizaje
automático en finanzas~\cite{kearns2013,lopezdeprado2018,bailey2014deflated}: una métrica de
clasificación puede ser muy alta y, aun así, el sistema no ser explotable si opera con
retraso, si aprende artefactos del propio conjunto de datos o si ignora el coste real de
cruzar una orden. Por eso el problema no se evalúa por el acierto, sino por el resultado
económico neto, una vez descontados costes y latencia realistas.

\section{Motivación}

La motivación de este trabajo es doble. Desde el punto de vista aplicado, los mercados de
predicción de criptoactivos son un entorno joven, con menos competencia algorítmica que
los mercados financieros tradicionales, y con datos públicos accesibles; resulta natural
preguntarse si contienen señales explotables. Desde el punto de vista metodológico —que es
el que da peso a una memoria de máster— el entorno es un excelente banco de pruebas para
una cuestión central en el aprendizaje automático aplicado a finanzas: la distancia entre
\emph{predecir correctamente} y \emph{obtener beneficio}.

¿Existe trabajo previo, y puede mejorarse? Existe, y con un hueco preciso: los estudios
académicos sobre Polymarket aparecidos en 2026 describen la plataforma —archivos de
eventos, arbitraje, composición del flujo— pero ninguno evalúa económicamente una
política de negociación bajo coste, latencia y comisiones medidos, ni valida
prospectivamente sus resultados (Capítulo~\ref{ch::estado}); y la literatura de
aprendizaje profundo sobre libros de órdenes rara vez cruza sus métricas de clasificación
con un modelo de ejecución realista. Ahí es donde este trabajo puede mejorar lo
existente, y ahí se juega su respuesta a «¿qué se gana?»: no una estrategia rentable
—el resultado económico final es negativo y se reporta íntegro—, sino la cuantificación,
por varias vías independientes, de cuánta de la capacidad predictiva sobrevive a la
ejecución, junto con el protocolo sellado que hace creíble esa cuantificación.

Esa distancia es el hilo conductor. Como se mostrará, un clasificador direccional sencillo
alcanza una precisión cercana al 90\,\% medida como \emph{markout} instantáneo y, pese a
ello, no sobrevive al primer escalón de latencia que el instrumento de captura permite
resolver. El coste y la ejecución, y no el acierto, determinan el resultado. Asumir lo
contrario es uno de los errores más comunes y costosos en este campo, y documentarlo con
evidencia cuantitativa es, en sí mismo, una aportación.

El argumento no descansa sobre una única medición. A lo largo del trabajo, la misma
conclusión se alcanza por tres caminos independientes: el efecto de la latencia sobre la
señal direccional, la estructura oficial de comisiones —que, una vez corregido un error
en el modelo de costes, resulta suficiente por sí sola para agotar el resultado bruto— y
la ausencia de cualquier retardo explotable en la incorporación de la información externa.
Que tres vías distintas converjan es lo que convierte una observación en una conclusión.

Conviene además situar las expectativas. El objetivo no es construir un sistema de trading en
producción ni demostrar rentabilidad, sino algo más modesto y, a la vez, más exigente
metodológicamente: determinar \emph{si} existe señal explotable bajo condiciones realistas y,
sobre todo, \emph{cómo} evaluarla sin engañarse. En un dominio donde es fácil obtener
resultados espectaculares pero ilusorios, el valor de un trabajo de máster está tanto en lo
que encuentra como en el rigor con el que descarta lo que no se sostiene. Por eso, a lo largo
de la memoria, los resultados negativos —modelos que parecían prometedores y no superaron la
validación— se reportan con el mismo cuidado que los positivos.

\section{Objetivos}

El objetivo general es \textbf{construir y validar, con un protocolo honesto, una
metodología completa para evaluar si el aprendizaje automático y profundo pueden anticipar
el movimiento de corto plazo del precio de los contratos de Bitcoin de Polymarket bajo
costes y latencia realistas}. De él se derivan los siguientes objetivos específicos:

\begin{enumerate}[noitemsep]
  \item Diseñar un sistema de captura de datos multifuente, con control de calidad por
  sesión, que produzca un corpus alineado temporalmente y libre de fugas.
  \item Definir un protocolo de evaluación que combine validación estrictamente temporal,
  un modelo explícito de costes y una latencia de ejecución medida contra el mercado real,
  declarando con precisión qué parte de ella es medición y qué parte es la resolución del
  instrumento.
  \item Cuantificar el impacto de la latencia sobre una señal direccional fuerte, para
  delimitar dónde se rompe la cadena entre predicción y ejecución.
  \item Evaluar una progresión de modelos, de menor a mayor complejidad (líneas base
  tabulares, modelos sensibles a la latencia y arquitecturas profundas sobre el libro de
  órdenes), exigiendo que cada incremento de complejidad se justifique con evidencia.
  \item Determinar si existe un candidato que sobreviva a una validación fuera de muestra
  bajo coste, y delimitar con rigor sus límites y su incertidumbre.
  \item Auditar el propio modelo de costes hasta sus fuentes primarias y, si resultara
  erróneo, corregirlo, documentar la corrección con marca de tiempo verificable y
  reevaluar sin reentrenar nada, extrayendo las consecuencias sobre qué lado del libro es
  económicamente viable.
  \item Someter la política resultante a una \textbf{validación prospectiva
  pre-registrada}: fijar por escrito y sellar criptográficamente, antes de observar los
  datos, el universo, los brazos, las métricas y las reglas de decisión; ejecutar una
  única evaluación sobre una ventana posterior al sellado; y reportar el resultado con
  independencia de su signo.
\end{enumerate}

Los dos últimos objetivos no figuraban en el planteamiento inicial y se incorporaron
durante el desarrollo: el sexto, al detectarse que el modelo de comisiones empleado no
correspondía con la tarifa oficial vigente; el séptimo, como consecuencia natural de que
toda la evidencia acumulada hasta ese punto era retrospectiva. Se declaran aquí como
objetivos, y no como incidencias, porque acabaron aportando más al trabajo que varios de
los previstos desde el principio.

\paragraph{Estructura del documento.} El presente trabajo se organiza como sigue:
\begin{itemize}[noitemsep]
  \item \textbf{Capítulo~\ref{ch::estado}} — Estado de la cuestión: literatura de
  microestructura y aprendizaje profundo sobre libros de órdenes, y sus limitaciones.
  \item \textbf{Capítulo~\ref{ch::metodos}} — Material y métodos: el sistema de captura,
  el corpus y su análisis exploratorio, las líneas base y la metodología de modelado,
  validación y costes.
  \item \textbf{Capítulo~\ref{ch::resultados}} — Resultados y análisis: la evaluación
  fuera de muestra, la fe de erratas de comisiones, el giro al lado \emph{maker} y la
  validación prospectiva pre-registrada.
  \item \textbf{Capítulo~\ref{ch::despliegue}} — Despliegue: arquitectura del sistema y
  condiciones bajo las que tendría sentido operarlo.
  \item \textbf{Capítulo~\ref{ch::etica}} — Aspectos éticos y limitaciones.
  \item \textbf{Capítulo~\ref{ch::reproducibilidad}} — Reproducibilidad: qué se publica
  y cómo reproducir las cifras.
  \item \textbf{Capítulo~\ref{ch::conclusiones}} — Conclusiones y trabajo futuro.
  \item \textbf{Anexos} — detalle de características e hiperparámetros, el pre-registro
  completo con el procedimiento de verificación del sello temporal, y una auditoría del
  registro público de operaciones desde el dato anonimizado.
\end{itemize}
En los preliminares se incluye además un glosario para el lector no familiarizado con las
criptomonedas.

Una advertencia sobre la lectura. La memoria no describe un plan ejecutado sin sobresaltos,
porque el trabajo no fue así. Recoge un itinerario en el que una señal aparentemente
excelente resultó no ser explotable, un modelo de costes tuvo que corregirse a mitad de
camino, el vehículo económico se desplazó de un lado del libro al contrario, y la validación
final —pre-registrada y sellada— arrojó un resultado negativo que se reporta íntegro. Ese
itinerario, y no un modelo concreto, es la aportación.
